\hypertarget{appendix-3f7e15effba6}{}\label{referencestart:appendix-3f7e15effba6}
\section{Named members versus array elements}\label{appendix-3f7e15effba6-named-members-versus-array-elements}

\hypertarget{appendix-3f7e15effba6-named-members-versus-array-elements}{}

\textbf{Curriculum depth:} Developed · \textbf{First seen:} C007 · \textbf{Latest development:} C008

\begin{allmightroledefinition}

\subsection{Definition and mechanism}\label{appendix-3f7e15effba6-definition-and-mechanism}

\hypertarget{appendix-3f7e15effba6-definition-and-mechanism}{}

The x member's one-past pointer is not an authorized y access. A span describes an already-valid range; adjacency cannot create that range.

\end{allmightroledefinition}

Adjacent x/y/z fields are separate subobjects, not T{[}3{]}. A pointer to x supports one-element-array pointer arithmetic; one-past cannot be dereferenced to access y. span(\&x,3) does not create an array. Use explicit member dispatch or actual array storage; a ToArray result is a copy.

\subsection{Prerequisites}\label{appendix-3f7e15effba6-prerequisites}

\hypertarget{appendix-3f7e15effba6-prerequisites}{}

\begin{itemize}
\tightlist
\item
  \href{../01-Bits-Programs-and-Reliable-Systems-P02/main.pdf\#nameddest=appendix-6af5921bc853}{Pointer arithmetic and provenance (Bits, Programs, and Reliable Systems, Part 2, Reference notes: Pointer arithmetic and provenance)} --- Array-relative pointer bounds explain why adjacency of separate members cannot authorize indexing.
\end{itemize}

\subsection{Used by}\label{appendix-3f7e15effba6-used-by}

\hypertarget{appendix-3f7e15effba6-used-by}{}

\begin{itemize}
\tightlist
\item
  \href{../07-Vision-XR-and-Playable-Worlds-P01/main.pdf\#nameddest=appendix-37a39984ba66}{Vector value types and component access (Vision, XR, and Playable Worlds, Part 1, Reference notes: Vector value types and component access)} --- Phase 6 uses this established foundation in its geometry narrative.
\end{itemize}

\subsection{Recall and next depth}\label{appendix-3f7e15effba6-recall-and-next-depth}

\hypertarget{appendix-3f7e15effba6-recall-and-next-depth}{}

Give an example, explain why the rule exists, and identify which guarantee would be missing if the prerequisite were ignored.

\textbf{Pending:} Deeper ownership, pointer/reference, generic programming and module rules from later sources. This is written coverage, not demonstrated personal mastery.

\subsection{Source and navigation}\label{appendix-3f7e15effba6-source-and-navigation}

\hypertarget{appendix-3f7e15effba6-source-and-navigation}{}

\hyperref[appendix-fd090d982a15-m008---parts-113-126]{C008 development and corrections}

\href{../01-Bits-Programs-and-Reliable-Systems-P02/main.pdf\#nameddest=01-foundations-i-ch003-bringing-the-object-model-back-to-bytes}{C007 teaching chapter (Bits, Programs, and Reliable Systems, Part 2, Ch. 3)} · \hyperref[appendix-fd090d982a15-m007---parts-101-112]{Source coverage} · \href{../07-Vision-XR-and-Playable-Worlds-P01/main.pdf\#nameddest=appendix-4ad47bdd269a}{Dependency Map (Vision, XR, and Playable Worlds, Part 1, Reference notes: Dependency Map)}
